\begin{frame}[t]{\ev{\tagM}Our API measurements cover three different workflows}
\vspace{.1cm}
{\small
\begin{tabularx}{\textwidth}{@{}Y r r r@{}}
\toprule
Workflow & Success & p50 & p95\\
\midrule
Daytona: create & 1,024/1,024 & 0.20 s & 1.12 s\\
Tensorlake: create & 256/256 & 3.44 s & 9.00 s\\
Our sandbox platform: restore, run, capture & 255/256 & 6.87 s & 9.04 s\\
\bottomrule
\end{tabularx}\par}
\vspace{.2cm}
{\small
\begin{ticks}
\item Restore path: a 141 MB named snapshot at concurrency 12
\item Whether the snapshot captures memory is unverified, so the third row is restore fan-out
\item Per-operation teardown is excluded from the percentiles
\item[\no] different operations and workloads: API round trips, not a ranking and not a memory-fork latency
\end{ticks}\par}
\claim{The table establishes exercised API paths and their latencies.}
\sources{Eliaz, \href{https://arxiv.org/abs/2607.09689}{arXiv:2607.09689}, Table 2.}
\note{[0:40] We measured three workflows. Daytona created 1,024 sandboxes, with a median of 0.20 seconds. Tensorlake created 256. On our sandbox platform, 255 of 256 round trips restored a 141 megabyte snapshot, ran a command and captured. The median was 6.87 seconds. The workloads differ, so the table is not a ranking. Memory capture is unverified. Next, we separate the calls the runtime must provide.}
\end{frame}
